
\paragraph{\textbf{Corollary 3.20.}}
\SGARefTarget{sga3:xiv:3.20}
{\itshape Let \(S\) be a prescheme and let \(G\) be a smooth
\(S\)-prescheme in groups of finite presentation over \(S\), whose
geometric fibers are connected reductive algebraic groups, that is,
extensions of a semisimple group by a torus. For every \(s\in S\)
there is an open neighborhood \(U\) of \(s\) such that
\(G|_U\) admits a maximal torus.}\footnote{This result, as well as
\SGARefLink{sga3:xiv:2.11}{2.11} on which it rests, immediately generalizes to the case in
which \(s\) is replaced by a finite subset of \(S\) contained in
an affine open.}

We shall see in Exposé XVI that the hypothesis on \(G\) implies
that \(G\) is affine over \(S\) and has locally constant reductive
rank. Hence, by \SGARefLink{sga3:xii:4.7}{XII~4.7}\textup{\SGARefLink{sga3:xii:4.7:item:c}{(c)}}, \(G\) admits a reductive center
\(Z\), the quotient
\[
  G'=G/Z
\]
is a smooth affine group over \(S\) whose reductive center is the
identity group, and the maximal tori of \(G\) correspond
bijectively to those of \(G'\). Moreover, the geometric fibers of
\(G'\) are plainly connected semisimple groups, and are adjoint by
definition because their reductive center is trivial.

It is therefore enough to treat the case where \(G\) is
semisimple and adjoint. By \SGARefLink{sga3:xiv:3.18}{3.18}, the problem reduces to finding an
open neighborhood \(U\) of \(s\) and a subgroup of type
\((\mathrm C)\) of \(G|_U\), or equivalently by \SGARefLink{sga3:xiv:3.9}{3.9}\textup{(a)}, a
Cartan subalgebra of \(\mathfrak g|_U\). This is possible when
\(k(s)\) is infinite: by \SGARefLink{sga3:xiv:3.7}{3.7}, \(\mathfrak g\) satisfies condition
\((\mathrm C_2)\) of \SGARefLink{sga3:xiv:2.9}{2.9}, and \SGARefLink{sga3:xiv:2.11}{2.11}\textup{(b)} applies.

In fact, \SGARefLink{sga3:xiv:3.20}{3.20} remains true without assuming that \(k(s)\) is
infinite. By the preceding argument, it is enough to know that, for
every adjoint semisimple group \(G\) over a finite field \(k\), the
Lie algebra \(\mathfrak g\) contains a regular element. Chevalley
proved this assertion using the properties of a Coxeter element of
the Weyl group; see the appendix below by J.-P. Serre.

\paragraph{\textbf{Remarks 3.21.}}
\SGARefTarget{sga3:xiv:3.21}
\textup{a)}
Assertion \SGARefLink{sga3:xiv:3.20}{3.20} remains valid, with essentially the same proof, if
\(G\) is replaced by a closed subprescheme in groups \(H\), smooth
over \(S\), whose rank is everywhere the same as that of \(G\)
\textup{(for example, a ``parabolic subgroup'' of \(G\))}, provided
that \(k(s)\) is infinite. I do not know whether the infinitude
hypothesis is again superfluous. It can at least be removed when
\(H\) is parabolic: constructing the radical \(U\) of \(H\) and
the quotient \(H/U\) reduces the question to the semisimple case.
Unfortunately, the method of regular elements seems powerless
here. One readily constructs examples---for instance, the
projective group with its ``standard'' Borel subgroup---in which
the Lie algebra \(\mathfrak h\) of \(H\) contains no regular
element.

\smallskip
\noindent
\textup{b)}
When \(k(s)\) is infinite, the proof of \SGARefLink{sga3:xiv:3.20}{3.20} in fact gives the
following sharper statement, using \SGARefLink{sga3:xiv:3.9}{3.9}\textup{(b)}: every maximal
torus \(T_0\) of \(G_s\) is induced by a maximal torus \(T\) over
an open neighborhood of \(s\). I do not know whether this remains
true without the infinitude hypothesis. The evident difficulty is
that the Lie algebra \(\mathfrak t_0\) of \(T_0\) need not contain
a regular element of the Lie algebra \(\mathfrak g_0\) of the
fiber \(G_s\).

An affirmative answer would imply the following assertion, which
is known only when the residue field is infinite or when \(A\) is
separated and complete. Let \(A\) be a local ring with residue
field \(k\), and let \(M\) be an ``Azumaya algebra'' over \(A\),
meaning an algebra that is a finite free \(A\)-module and whose
reduction
\[
  M_0=M\otimes_A k
\]
is a central simple algebra over \(k\). Let \(D_0\subset M_0\) be
a commutative subalgebra separable over \(k\), with
\[
  [M_0:k]=[D_0:k]^2.
\]
Then there should exist a commutative subalgebra \(D\subset M\)
that is a direct summand as an \(A\)-module and satisfies
\[
  D\otimes_A k=D_0.
\]

The datum of \(M\) is equivalent to that of a principal homogeneous
bundle under the projective group \(\operatorname{PGL}(n)_A\), and
hence to an inner twisted form \(G\) of \(\operatorname{PGL}(n)\).
The maximal tori of \(G\) correspond bijectively to the commutative
subalgebras \(D\subset M\) that are étale of rank \(n\) over \(A\).

\smallskip
\noindent
\textup{c)}
Applying \SGARefLink{sga3:xiv:3.20}{3.20} to the centralizer of a subtorus \(Q\) of a reductive
\(S\)-prescheme in groups \(G\), one sees that every such \(Q\) is
contained, Zariski locally, in a maximal torus of \(G\).
